\section{Organización de la Memoria}

La memoria comienza con el \hyperref[cap:tecnologiasUtilizadas]{Capítulo~\ref*{cap:tecnologiasUtilizadas}. Tecnologías Utilizadas}, donde se presentan las principales herramientas, modelos y técnicas que sirven de base al sistema. En él se introducen, entre otros, los métodos de detección, tracking, análisis geométrico, modelado posicional, detección secuencial de eventos, generación de lenguaje y síntesis de voz.

A continuación, el \hyperref[cap:trabajoRelacionado]{Capítulo~\ref*{cap:trabajoRelacionado}. Trabajo Relacionado} revisa trabajos previos en los distintos bloques del pipeline. Esta revisión permite situar el proyecto dentro del estado del arte y contextualizar las decisiones adoptadas en tareas como la detección de jugadores y balón, el tracking multiobjeto en vídeo deportivo y la detección de eventos futbolísticos.

Una vez establecido ese contexto, el \hyperref[cap:procesamientoVisual]{Capítulo~\ref*{cap:procesamientoVisual}. Procesamiento Visual} describe la primera gran fase del sistema. En este capítulo se explica cómo se transforma el vídeo broadcast en una representación estructurada del partido mediante detección de elementos del juego, calibración geométrica del campo, clasificación de equipos, una primera capa de asociación temporal basada en ByteTrack, una segunda capa de tracking canónico y, finalmente, una evaluación heurística orientada a depurar y seleccionar la configuración final del pipeline visual.

Sobre esa salida se construye el \hyperref[cap:analisisTactico]{Capítulo~\ref*{cap:analisisTactico}. Análisis Táctico}, donde se abordan tareas de mayor nivel semántico. En él se describe, en primer lugar, una inferencia heurística de la posesión del balón. Después se presenta la inferencia de posiciones futbolísticas a partir de la disposición espacial de los jugadores y de su contexto colectivo. Finalmente, se expone la clasificación de acciones relevantes del partido, como pases, controles, robos o tiros.

La siguiente etapa se presenta en el \hyperref[cap:narracion]{Capítulo~\ref*{cap:narracion}. Narración}, centrado en la generación automática de comentarios. Este capítulo resume cómo los eventos detectados se convierten en texto narrativo mediante un modelo de lenguaje y cómo ese texto pasa posteriormente a audio mediante distintos sistemas de síntesis de voz evaluados durante el proyecto.

Antes del cierre, el \hyperref[cap:integraciónSistema]{Capítulo~\ref*{cap:integraciónSistema}. Integración del Sistema y Validación Final} explica como se consume el sistema desde una interfaz web y presenta una validación global del comportamiento del sistema en condiciones finales de uso.

Finalmente, la memoria concluye con el \hyperref[cap:conclusiones]{Capítulo~\ref*{cap:conclusiones}. Conclusiones y Trabajo Futuro}, donde se sintetizan los principales resultados obtenidos, se analizan las limitaciones del sistema y se recogen posibles líneas de mejora para continuar el desarrollo del trabajo.


